\slajdid{09_3-s054}
\begin{frame}[label=09_3-s054]
\gusttekst\setlength{\abovedisplayskip}{3pt}\setlength{\belowdisplayskip}{3pt}\setlength{\abovedisplayshortskip}{2pt}\setlength{\belowdisplayshortskip}{2pt}Ono što takođe treba da uočite da u 1. slučaju koji smo imali dolazi principski do još veće degradacije naponskih nivoa pošto je tada sors nMOS tranzistora na visokom potencijalu, različitom od potencijala osnove pa dolazi i do izražaja promene praga $V_{Tn}$ zbog različitih potencijala sorsa i osnove.
\par\medskip U tom smislu ako bi želeli da napravimo višeulaznu logiku na način
\begin{center}\begin{tikzpicture}[x=.85cm,y=.53cm,font=\sitneoznake,>=Latex]\ctikzset{bipoles/length=.6cm}\begin{scope}[shift={(0,0)}]\draw(-.35,0)--(.35,0);\draw(-.4,-.13)--(.4,-.13);\draw(-.3,0)--(-.3,.4)--(-1,.4)node[ocirc]{}node[left]{$B=V_{DD}$};\draw(.3,0)--(.3,.4)--(1,.4);\draw(0,-.13)--(0,-.8)node[ocirc]{}node[below]{$A=V_{DD}$};\node[above]at(0,.1){1};\end{scope}\begin{scope}[shift={(1,1.5)}]\draw(-.35,0)--(.35,0);\draw(-.4,-.13)--(.4,-.13);\draw(-.3,0)--(-.3,.4)--(-1,.4)node[ocirc]{}node[left]{$C=V_{DD}$};\draw(.3,0)--(.3,.4)--(1,.4);\draw(0,-.13)--(0,-.8)node[ocirc]{};\node[above]at(0,.1){2};\end{scope}\begin{scope}[shift={(2,3)}]\draw(-.35,0)--(.35,0);\draw(-.4,-.13)--(.4,-.13);\draw(-.3,0)--(-.3,.4)--(-1,.4)node[ocirc]{}node[left]{$D=V_{DD}$};\draw(.3,0)--(.3,.4)--(1,.4);\draw(0,-.13)--(0,-.8)node[ocirc]{};\node[above]at(0,.1){3};\end{scope}\draw(1,.4)--(1,.7);\node[right]at(1,.5){$V_{DD}-V_{Tn1}$};\draw(2,1.9)--(2,2.2);\node[right]at(2,2){$V_{DD}-V_{Tn1}-V_{Tn2}$};\draw(3,3.4)--(6.8,3.4);\node[above right]at(2.5,3.65){$Y=V_{DD}-V_{Tn1}-V_{Tn2}-V_{Tn3}$};\draw(6.4,3.4)to[C,l=$C_L$](6.4,1.8)node[ground]{};\end{tikzpicture}
\end{center}
u situaciji prikazanoj na slici došlo bi do značajne degradacije logičkih signala na izlazu pri čemu su zbog različitih polarizacija sorsa u odnosu na osnovu i različiti pragovi $V_{Tn}$ pojedinih tranzistora. Osnove su kao i uvek na potencijalu mase.
\end{frame}
